\documentclass{article}
% Damit die Verwendung der deutschen Sprache nicht ganz so umst\"andlich wird,
% sollte man die folgenden Pakete einbinden: 
\usepackage[latin1]{inputenc}% erm\"oglich die direkte Eingabe der Umlaute 
\usepackage[T1]{fontenc} % das Trennen der Umlaute
\usepackage{ngerman} % hiermit werden deutsche Bezeichnungen genutzt und 
                     % die W\"orter werden anhand der neue Rechtschreibung 
		     % automatisch getrennt.  
\usepackage{a4wide}            
\title{Synthetic Data generation techniques}
%\author{Ihr Name  \\
	%Ihr Unternehmen / Universit\"at  \\
	%Teststra\ss e -99 \\
	%0123456 Testhausen \\
	%\and 
	%Der Andere  \\
	%Sein Unternehmen / Universit\"at \\
	%Musterstra\ss e 00 \\
	%6543210 Musterdorf \\
	%}

\date{\today}
% Hinweis: \title{um was auch immer es geht}, \author{wer es auch immer 
% geschrieben hat} und  \date{wann auch immer das war} k\"onnen vor 
% oder nach dem  Kommando \begin{document} stehen 
% Aber der \maketitle Befehl mu\ss{} nach dem \begin{document} Kommando stehen! 
\begin{document}

\maketitle


\begin{abstract}
Eine kurze Zusammenfassung um was es in der vorliegenden Arbeit \"uberhaupt
so geht \ldots
\end{abstract}

\section{Einf\"uhrung}
% Was sind synthetische Daten ?
Ziel der Arbeit ist es m\"oglichsten vielen oder wenn m\"oglichen es allen 
zu erm\"oglichen, Dokumente mit \LaTeX{} zu erstellen!

\subsection{Why}
% Wozu braucht man synthetische Daten
- Beantragungsverfahren dauer lange und Anträge können abgelehnt werden
- Rapid Prototyping (Testen eines Dashboards oder Pipeline)
- Regulatorische Compliance -> Teilen von Daten
% siehe duck.ai am Laptop

\subsection{Challenges}
%Herausforderungen beim erstellen synthetischer Daten
- Qualität vs Privacy
- wiss. Veröffentlichung (qualität vs genug privacy)
- Tailoring (jedes Modell trainiert genau eine Aufgabe)
- Wie misst man die Qualität (welche Metriken)
- Welche Methodik für welchen Art an Daten

\section{How a CTGAN works}
Ziel der Arbeit ist es m\"oglichsten vielen oder wenn m\"oglichen es allen 
zu erm\"oglichen, Dokumente mit \LaTeX{} zu erstellen!

\subsection{Aufbau}
\subsection{Schwierigkeiten}
\subsection{Evaluierung}

\section{Introduction to table diffusion}
\subsection{Aufbau}
\subsection{Schwierigkeiten}
\subsection{Evaluierung}

\section{Comparison}

\section{Diskussion}

\section{Ausblick}
- TVAE
- LLM
- TAbddpm
- SimbaML (simulation based machine learning <-- ml augemented SIR equation fit)

\section{Quellen}

\paragraph{Modellkonzeption}
Zu Beginn fangen wir mit einem kleinen Beispiel f\"ur die article Klasse an,
diese ist eine sehr wichtige Dokumentenklasse. Aber daneben gibt es noch
weitere wie book \ref{book}, report \ref{report} und letter \ref{letter} 
welche im Abschnitt \ref{documentclasses} beschrieben werden. Letztendlich 
wird im Abschnitt \ref{conclusions} ein Fazit gezogen.



%\section{Dokumentenklassen} \label{documentclasses}

%\begin{itemize}
%\item article
%\item book 
%\item report 
%\item letter 
%\end{itemize}


%\begin{enumerate}
%\item article
%\item book 
%\item report 
%\item letter 
%\end{enumerate}

%\begin{description}
%\item[article\label{article}]{Article ist \ldots}
%\item[book\label{book}]{Die book Klasse ist \ldots}
%\item[report\label{report}]{Die Klasse report erm\"oglicht es  \ldots}
%\item[letter\label{letter}]{Wenn man einen Breif schreiben sollte man eine 
%	andere Klasse nutzen, da diese f\"ur ein anderes als das deutsche 
	%Briefformat ausgelegt ist.}
%\end{description}


%\section{Fazit}\label{conclusions}
%Nach langer Suche hat sich herausgestellt, dass es kein l\"angeres 
%\LaTeX{} Beispiel, als das von \cite{doe} geschriebene gibt. 

\begin{thebibliography}{9}
\bibitem[Doe]{doe} \emph{Erstes und letztes \LaTeX{} Beispiel.},
John Doe 50 v.Chr.  
\end{thebibliography}

\end{document}